% ============================================
% MEM-08
% Gerado de molambudos.md — build_miolo.py
% ============================================
\section*{\textbf{MEM-08} --- A Chegada ao Colônia}
\markboth{MEM-08 --- A Chegada ao Colônia}{}

\noindent \textit{1917 --- Barbacena, Minas Gerais}

\noindent \textbf{1917 --- Barbacena, Minas Gerais}

\vspace{0.5em}\noindent\rule{\textwidth}{0.5pt}\vspace{0.5em}

\noindent Eu não lembro como cheguei no Colônia.

\noindent Lembro do caminhão. Lembro de sacos. Lembro que estava escuro e que a gente era muitos dentro de um espaço pequeno. Gente que cheirava a suor, a fezes, a medo.

\noindent Depois o caminhão parou.

\noindent Eu vi o portão.

\noindent Era um portão de ferro, preto, com pontas em cima. Atrás do portão tinha uma estrada de terra. Depois da estrada tinha um prédio. O prédio era grande, cinza, com janelas pequenas e uma cruz no topo.

\noindent Cruz de hospital. Mas parecia cruz de cemitério.

\noindent \textit{\textquotedblleft{}Desce\textquotedblright{}}, alguém disse.

\noindent Eu desci.

\noindent \textit{\textquotedblleft{}Nome?\textquotedblright{}}

\noindent \textit{\textquotedblleft{}Joaquim Antônio Correia.\textquotedblright{}}

\noindent \textit{\textquotedblleft{}Idade?\textquotedblright{}}

\noindent \textit{\textquotedblleft{}Dez.\textquotedblright{}}

\noindent O homem olhou pra mim. Olhou pro meu olho amarelo. Anotou alguma coisa num papel.

\noindent \textit{\textquotedblleft{}Diagnóstico?\textquotedblright{}} ele perguntou pra outro homem.

\noindent \textit{\textquotedblleft{}Veio do Ceará. Curral do governo. Órfão. Não fala.\textquotedblright{}}

\noindent \textit{\textquotedblleft{}Eu falo\textquotedblright{}}, eu disse.

\noindent Os dois homens me olharam.

\noindent \textit{\textquotedblleft{}Então fala\textquotedblright{}\textit{, disse o primeiro. }\textquotedblleft{}O que tu tem?\textquotedblright{}}

\noindent Eu pensei em dizer a verdade. Pensei em falar da vala. Da criatura. Do olho. Da fome que nunca passa.

\noindent Mas eu tinha aprendido uma coisa na estrada e no curral: a verdade não ajuda.

\noindent \textit{\textquotedblleft{}Nada\textquotedblright{}\textit{, eu disse. }\textquotedblleft{}Não tenho nada.\textquotedblright{}}

\noindent O homem riu.

\noindent \textit{\textquotedblleft{}Todo mundo que chega aqui tem alguma coisa. Se não tem, a gente inventa.\textquotedblright{}}

\noindent Ele escreveu no papel:

\noindent \textit{\textquotedblleft{}Paciente 1.260 --- Joaquim Antônio Correia. Esquizofrenia hebefrênica. Perigoso. Isolamento.\textquotedblright{}}

\noindent Eu não sabia o que era esquizofrenia. Aprendi depois. É uma palavra que os médicos usam quando não sabem o que uma pessoa tem.

\noindent Eles me levaram para uma sala.

\noindent A sala tinha uma cama de ferro. Uma janela gradeada. Um balde no canto.

\noindent \textit{\textquotedblleft{}Fica aqui\textquotedblright{}\textit{, disse o enfermeiro. }\textquotedblleft{}Amanhã vem o médico.\textquotedblright{}}

\noindent \textit{\textquotedblleft{}E a comida?\textquotedblright{}}

\noindent \textit{\textquotedblleft{}Amanhã.\textquotedblright{}}

\noindent Ele fechou a porta. A chave girou.

\noindent Eu fiquei ali, sentado na cama, olhando para a parede. A parede estava cheia de riscos. Riscos de unha. Riscos de dente. Riscos de gente que tinha ficado ali antes de mim.

\noindent \textit{\textquotedblleft{}Bem-vindo\textquotedblright{}}, disse a criatura dentro de mim.

\noindent \textit{\textquotedblleft{}Cala a boca.\textquotedblright{}}

\noindent \textit{\textquotedblleft{}Eu vou gostar daqui\textquotedblright{}\textit{, ela disse. }\textquotedblleft{}Aqui tem muita fome.\textquotedblright{}}

\noindent Ela estava certa.

\noindent O Colônia tinha mais de cinco mil pacientes. Os dormitórios eram galpões com camas encostadas umas nas outras. A comida era um caldo de osso com farinha. A água era servida uma vez por dia.

\noindent Os gritos não paravam nunca.

\noindent De noite, os gritos ecoavam pelos corredores. Gritos de dor. Gritos de medo. Gritos de gente que não sabia mais quem era.

\noindent Nos primeiros dias, eu tampava os ouvidos.

\noindent Depois de uma semana, eu já gritava junto.

\noindent A criatura crescia.

\noindent Cada grito alimentava ela. Cada lágrima. Cada pedaço de comida que eu engolia e que não saciava.

\noindent Ela se alimentava do Colônia.

\noindent E o Colônia se alimentava de mim.

\vspace{0.5em}\noindent\rule{\textwidth}{0.5pt}\vspace{0.5em}

\noindent Anos depois, Dr. Oliveira me perguntou:

\noindent \textit{\textquotedblleft{}Joaquim, você sabe quantos anos passou no Colônia?\textquotedblright{}}

\noindent Eu sabia.

\noindent \textit{\textquotedblleft{}Sessenta e dois.\textquotedblright{}}

\noindent \textit{\textquotedblleft{}E como você se sente sobre isso?\textquotedblright{}}

\noindent Eu olhei pra ele. O olho amarelo brilhou.

\noindent \textit{\textquotedblleft{}Eu não vivi aqui, doutor. Eu apodreci aqui.\textquotedblright{}}

\noindent Ele anotou.

\noindent \textit{\textquotedblleft{}Apodrecer\textquotedblright{}\textit{, ele repetiu. }\textquotedblleft{}Interessante escolha de palavras.\textquotedblright{}}

\noindent \textit{\textquotedblleft{}Não foi escolha\textquotedblright{}\textit{, eu disse. }\textquotedblleft{}Foi o que aconteceu. A gente não escolhe apodrecer. A gente só... apodrece.\textquotedblright{}}

\noindent Ele parou de anotar.

\noindent \textit{\textquotedblleft{}O que aconteceu com você, Joaquim?\textquotedblright{}}

\noindent Eu sorri.

\noindent \textit{\textquotedblleft{}O que não aconteceu?\textquotedblright{}}

\vspace{0.5em}
\noindent\textit{\textbf{↪ Rotas:} \rota{MEM-09} \rotasep \rota{DOC-02} \rotasep \rota{MEM-12}}